\section{Multiplicity comparison for the native isotypic decomposition}

% BEGIN F86I1
\paragraph{F86I1. The parameter of a type already present at the reference level.}
Fix $\sigma$ with $n_0(\sigma)>0$. Put $G_Y=U(Y_0)$ and
$G_W=U(W_j)$ for $1\le j\le m$. We first work on $G_Y$ alone.
A reference basis row gives a nonzero vector in the level part of the
$\bar\sigma$-isotypic cuspidal multiplicity space. Complex conjugation
of both the functions and the archimedean model identifies it
antilinearly with the $\sigma$-part for the conjugate discrete series
$M_0$. Its level character is now the one read on the lift side:
the level values have modulus one. Conjugation preserves the good
projector. Indeed G85M3 pairs each finite series with its
reciprocal; the good set is closed under this operation, so their
sum of central idempotents is fixed. We thus have an actual
cuspidal automorphic representation on $G_Y$ with finite part
$\tau_f$ representing $\sigma$ and the archimedean factor prescribed
by $M_0$. This follows by decomposing the cuspidal spectrum at the
given finite level and discrete series. In the convention using
multiplicity \emph{maps} it is the archimedean representation
intertwined by those maps. The same convention is used on $G_W$.

The two real groups can have different signatures. The reference
group has signature $(k+r_S,r_S+1)$, and the $j$-th target has
signature $(k+j-1,m-j)$. Their specified real discrete series
have the same unordered Harish--Chandra parameters at each
embedding: these are the $r$ entries $\lambda_a$ of the member's
$\mathrm{lamAt}$, arranged according to the positive and negative
slots of the respective signature. In the reference model the
preceding conjugation is essential for this equality. The distinguished
place on the lift reads the original archimedean labels. The entries are integral. Here one takes the intersection with the
eventual host range for the frame and the weights, before choosing
the member. In that range $h,\Delta\ge256N(m)^4$ and the first
node $b_0\ge100r^4$. The consecutive gaps in the upper segment of
the \emph{rank-$r$} tuple are $4\Delta$, the middle gap is $b_0$,
and the later gaps are $b_{j+1}-b_j=2h$ by the node and endpoint
equations. Since $r\le N(m)$ these gaps exceed one, uniformly
in the original row index. This uses the maximum of the original
size and first-node thresholds, as in the restriction of family thresholds, and gives
a strongly regular infinitesimal character on both groups.

Use Labesse, \emph{Changement de base CM et s\'eries discr\`etes},
Corollary~5.3, on this reference representation. The CM extension is
unramified at its finite nonsplit places and $[F:\mathbb Q]>1$ as
in A72F6 (indeed its degree is even by the quadratic class
character evaluated at $-1$). Labesse's sign hypothesis $(*)$
also holds for this possibly noncompact real group. For detail,
the strongly regular case of Salamanca--Riba's unitary dual theorem
(Duke Math.\ J.\ 96 (1999), 521--546) expresses a unitary module
at this infinitesimal character as cohomological induction from a
unitary character of a $\theta$-stable Levi. The Levi blocks of
$U(p,q)$ are unitary groups. In a block of size $d$ the inducing
character is constant on its $d$ coordinates and the block's
half-sum of positive roots contributes a string with step one;
the contribution from the nilradical is constant on the block.
A block of size at least two would give two entries differing by
one in the infinitesimal character. Consequently the inducing
Levi is a compact Cartan and the resulting modules are discrete
series in its regular real packet. This is the sufficient criterion
for $(*)$ given immediately before Labesse's corollary.

The corollary gives an isobaric sum of conjugate-self-dual
discrete representations of general linear groups, with local
base-change equality at \emph{all split places} and at almost
every nonsplit place. Write the discrete factors as Speh
representations of cuspidal bases. At a designated type place
the inertia polynomial on wild invariants has degree $r$ and is
squarefree, by G85M1--3. Complex conjugation of the original
type, if used to read the polynomial, reciprocates it and
preserves both tests. The degree of a factor on wild invariants
is bounded by its rank, and the sum of all ranks with lengths
is $r$. Thus every factor has full degree there. Norm shifts
are trivial on inertia, so a length $d>1$ would repeat every
root of a nonconstant factor $d$ times. All lengths are one.
If there were more than one base, a base of smallest rank
$1\le e\le r/2$ would give a degree-$e$ divisor at every designated place.
The divisor is defined over the residue field: Frobenius permutes
the tame eigenvalues by the residue-field power, and it preserves
each Weil summand. The good row has a place admitting no such
divisor for each $1\le e\le r/2$. This proves that the base change is
one rank-$r$ cuspidal representation $\Pi$. The cuspidal clause
at the end of Labesse's corollary then gives the actual
archimedean base change of the reference discrete series, so the
real packet used below is the one for $\Pi$. The argument started
with $n_0(\sigma)>0$ and makes no assumption about the occurrence
of $\sigma$ on $W_j$.
% END F86I1

% BEGIN F86I2
\paragraph{F86I2. Multiplicity one on both inner forms and the original field \textnormal{\texttt{adm8}}.}
Use the generic pure-inner-form case of Kaletha--M\'inguez--Shin--White,
\emph{Endoscopic classification of representations: inner forms
of unitary groups}, arXiv:1409.3731, Corollary~3.3.2 and
Theorem~1.7.1 (the unconditional case specified directly after the
theorem). Both algebraic groups arise from Hermitian spaces:
an isometry over $\overline F$ to the quasi-split space gives a
Galois cocycle with values in its unitary group. Thus these are
pure inner twists. The spectral support theorem assigns an Arthur
parameter to the near-equivalence class of the reference
representation. Take the ordinary stable base-change normalization
for odd rank. The general-linear class of this parameter is determined
by its unramified transfers at almost all finite places. Labesse's
unramified comparisons in F86I1 give precisely the same class as
$\Pi$; this includes the inert as well as the split places. Uniqueness
of the expanded isobaric list from almost-everywhere agreement
(A72F5 or S86A1) identifies the parameter as the single cuspidal
$\Pi$ of Arthur length one. In particular it is generic in the
sense required in the cited theorem. Its reduced global component
group is trivial: the centralizer of a single multiplicity-one
self-dual summand consists of the two scalar signs, and both
belong to the Galois-fixed center of the dual group.
Every tensor of local packet members on either pure inner form
therefore has discrete multiplicity one.

Take on $G_W$ the finite factors of $\tau_f$ transported by
the member's finite-adelic isometry. Its conjugation of the local
groups preserves the local packet; changing the frame by an inner
automorphism gives an isomorphic representation. At infinity
choose the stipulated discrete-series model for $W_j$. The real
discrete-series correspondence for unitary groups parametrizes the
packet across signatures by the allocations of a fixed ordered
regular tuple to the positive and negative slots. F86I1 verifies
that this very model has the reference tuple, at every real
place. This proves occurrence with multiplicity one. The resulting
discrete representation is cuspidal. Indeed otherwise its
cuspidal support would be on a proper rational Levi. A proper
parabolic of a unitary group stabilizes an isotropic flag; the
dimension $e$ of its first nonzero proper step is fixed over $F$.
At every designated split place this creates a general-linear
Levi factor of rank $e$, and the local constituent is a subquotient
of induction from that Levi. The supercuspidal support of a
subquotient of such induction is the union of the induced
supports. Unramified changes in these supports do not change
their tame eigenvalues. Full wild-invariant degree therefore gives
a residue-field polynomial divisor of degree $e$ at every such
place, the same contradiction as in F86I1. This argument also
shows why residual discrete representations cannot add a copy
to the cuspidal spaces in this comparison.

It remains to compare the dimension cut out in one finite copy
on each side. The finite isometry can be chosen to match the two
lattice frames at the level primes: their integral Gram frames
agree there and the comparison has identity reduction at the
specified type primes. It carries the base subgroup and its
character to the lift subgroup and its character. On finite
vectors the other cut is the same finite sum of right translations
by lifts of the residue group, with the same good central
coefficients. Replacing a lift by another changes it by the
principal congruence kernel, which acts trivially here. The
frame used to name $\tau_f$ differs from this isometry by inner
conjugation. Thus the two cut spaces within the fixed finite
representation are isomorphic, with any conjugation of the
original reference reading already included in F86I1. The norm
factor $|P_0|$ from the realization compares squared Hilbert
norms and leaves their dimensions unchanged.

The reference $\sigma$-cut consequently has the dimension of
one copy of this finite cut. Explicitly the vectors with that
label from the orthonormal reference basis give $n_0(\sigma)$
independent vectors in it. The cut in $M_j$ contains as many
vectors, since the transported copy with the prescribed infinite
type has just been exhibited. If $n_0(\sigma)=0$ the required
lower bound is automatic. We have proved the bound in L86M5 for
every $\sigma$. Its dimension sum and independence now give
$\displaystyle M_j=\bigoplus_\sigma\mathrm{isot}_j(\sigma)$,
and in fact $\dim\mathrm{isot}_j(\sigma)=n_0(\sigma)$ for every
$\sigma$. This is the original $\mathrm{adm8}$ for all
$1\le j\le m$. Index $0$ is the reference row providing $n_0$.
% END F86I2
